% config/typography.tex — Fonts, math setup, micro-typography
% ------------------------------------------------------------
% Requires LuaLaTeX (fontspec + unicode-math, loaded in main.tex
% BEFORE this file). Loaded after config/colors.tex.

% ---- Font set toggle ----------------------------------------
% pagella : TeX Gyre Pagella + Pagella Math  (warm "olympiad book" look)
% modern  : Latin Modern + LM Math            (classic journal look)
\newif\ifpagellafonts
\pagellafontstrue   % <- flip to \pagellafontsfalse for Latin Modern

\ifpagellafonts
  \setmainfont{TeX Gyre Pagella}[
    Ligatures   = TeX,
    BoldFont    = *-Bold,
    ItalicFont  = *-Italic,
    BoldItalicFont = *-BoldItalic,
  ]
  \setsansfont{TeX Gyre Heros}[
    Scale       = 0.92,
    BoldFont    = *-Bold,
    ItalicFont  = *-Italic,
    BoldItalicFont = *-BoldItalic,
  ]
  \setmonofont{Latin Modern Mono}[Scale=0.88]
  \setmathfont{TeX Gyre Pagella Math}
\else
  \setmainfont{Latin Modern Roman}
  \setsansfont{Latin Modern Sans}[Scale=0.92]
  \setmonofont{Latin Modern Mono}[Scale=0.88]
  \setmathfont{Latin Modern Math}
\fi

% Sans font is reserved for labels, running heads, box titles and
% difficulty badges; body and math stay serif for readability.
\newcommand{\labeltext}[1]{{\sffamily\bfseries\color{navy}#1}}

% ---- Micro-typography ---------------------------------------
% Protrusion + font expansion: subtle but lifts the page a full notch.
\usepackage{microtype}
\microtypesetup{
  protrusion = true,
  expansion  = true,
  factor     = 900,
}

% ---- Spacing conventions ------------------------------------
% A touch more leading than default: long math-heavy paragraphs
% with display equations read better.
\usepackage{setspace}
\setstretch{1.04}

% Display math: slightly tighter spacing around equations than default.
\AtBeginDocument{%
  \setlength{\abovedisplayskip}{6pt plus 2pt minus 3pt}%
  \setlength{\belowdisplayskip}{6pt plus 2pt minus 3pt}%
  \setlength{\abovedisplayshortskip}{3pt plus 1pt}%
  \setlength{\belowdisplayshortskip}{3pt plus 1pt}%
}

% Never let a paragraph start on the same line as a display above it.
\usepackage{parskip}%  -- deliberately NOT activated; this template uses
                     %  classic indentation for dense mathematical prose.
% Paragraph indent + skip tuned for curriculum prose:
\setlength{\parindent}{1.15em}
\setlength{\parskip}{2pt plus 1pt}
